
\documentclass[11pt,a4paper]{article}

\usepackage[margin=1in]{geometry}
\usepackage{amsmath,amssymb,mathtools,bm}
\usepackage{booktabs,array,longtable,multirow}
\usepackage{enumitem}
\usepackage{xcolor}
\usepackage{hyperref}
\usepackage{algorithm}
\usepackage{algpseudocode}
\usepackage{graphicx}
\usepackage{microtype}
\usepackage{setspace}
\usepackage{titlesec}

\hypersetup{
    colorlinks=true,
    linkcolor=blue!60!black,
    urlcolor=blue!60!black,
    citecolor=blue!60!black
}

\setlength{\parskip}{0.45em}
\setlength{\parindent}{1.6em}
\setstretch{1.12}

\titleformat{\section}{\Large\bfseries\color{blue!50!black}}{\thesection}{0.6em}{}
\titleformat{\subsection}{\large\bfseries\color{blue!35!black}}{\thesubsection}{0.6em}{}
\titleformat{\subsubsection}{\normalsize\bfseries\color{blue!30!black}}{\thesubsubsection}{0.6em}{}

\newcommand{\E}{\mathbb{E}}
\newcommand{\Var}{\operatorname{Var}}
\newcommand{\Cov}{\operatorname{Cov}}
\newcommand{\norm}[1]{\left\lVert #1 \right\rVert_2}
\newcommand{\clip}{\operatorname{clip}}
\newcommand{\CV}{\operatorname{CV}}
\newcommand{\Corr}{\operatorname{Corr}}
\newcommand{\sgn}{\operatorname{sgn}}

\title{\textbf{From Effective Signal Quality in GRPO to Learning-Rate Control:\\
RL Noise Decomposition, Empirical Findings, and Algorithmic Directions}}
\author{}
\date{May 20, 2026}

\begin{document}
\maketitle

\begin{abstract}
This report summarizes our updated analysis of learning-rate sensitivity in GRPO/PPO-style LLM reinforcement learning. The original goal was to derive an optimal learning-rate form from \(L\)-smoothness:
\[
\alpha_t^\star
=
\frac{1}{L_t}
\cdot
\frac{\|g_t\|^2}{\|g_t\|^2+N_t},
\]
and then estimate the signal fraction. However, our experiments show that both batch-split \(r_t\) and momentum-based estimators are unreliable in the current LLM RL setting. Stage-2 and independent \(A^2,Q,H\) analyses suggest a different core issue: the effective policy-gradient signal is sparse, heavy-tailed, and concentrated on a small subset of prompts and responses. In particular, the response-level policy-gradient energy proxy
\[
H_{i,k}=A_{i,k}^2Q_{i,k}
\]
is highly concentrated, and the variance of \(H\) is primarily driven by the advantage-energy term \(A^2\) and its interaction with score/length energy \(Q\). Therefore, sampling-side correction is the more direct root-cause intervention, while LR scheduling is better understood as a low-intrusion risk-control mechanism. The recommended next direction is to move away from the previous PG-Moment/\(\lambda_M\) calibration scheme and toward \textbf{Signal-Quality-Aware LR Scaling}, where the base learning rate is modulated by effective advantage-signal availability and gradient-energy concentration risk.
\end{abstract}

\tableofcontents
\newpage

\section{Background: From ``Best LR'' to RL Signal Quality}

\subsection{Optimal LR from \(L\)-Smoothness}

Consider maximizing an RL objective:
\[
J(\theta)
=
\E_{x\sim \mathcal D,\ y\sim \pi_\theta(\cdot|x)}
\left[R(x,y)\right].
\]
Assume that \(J\) is locally \(L_t\)-smooth around \(\theta_t\). With a stochastic policy-gradient update
\[
\theta_{t+1}=\theta_t+\alpha \hat g_t,
\]
we have the lower bound
\[
J(\theta_t+\alpha \hat g_t)
\ge
J(\theta_t)
+
\alpha g_t^\top \hat g_t
-
\frac{L_t}{2}\alpha^2\|\hat g_t\|^2,
\]
where
\[
g_t=\nabla J(\theta_t).
\]
If
\[
\hat g_t = g_t+\epsilon_t,\qquad \E[\epsilon_t]=0,
\]
and
\[
N_t=\E\|\epsilon_t\|^2,
\]
then
\[
\E\|\hat g_t\|^2=\|g_t\|^2+N_t.
\]
The expected lower-bound improvement is therefore
\[
\E[\Delta J_{lb}(\alpha)]
=
\alpha \|g_t\|^2
-
\frac{L_t}{2}\alpha^2(\|g_t\|^2+N_t).
\]
Maximizing this quadratic in \(\alpha\) gives
\[
\alpha_t^\star
=
\frac{1}{L_t}
\cdot
\frac{\|g_t\|^2}{\|g_t\|^2+N_t}.
\]
Define the signal fraction
\[
r_t
=
\frac{\|g_t\|^2}{\|g_t\|^2+N_t}.
\]
Then
\[
\alpha_t^\star=\frac{r_t}{L_t}.
\]

This decomposition gives two important messages:
\begin{enumerate}[leftmargin=2em]
    \item The optimal LR should decrease when the gradient signal fraction decreases.
    \item Even if \(r_t\) were estimated well, the absolute scale \(1/L_t\) remains. Therefore, a first practical method should be positioned as LR modulation:
    \[
    \alpha_t=\alpha_{\mathrm{base}} q_t,
    \]
    where \(\alpha_{\mathrm{base}}\) captures the average curvature scale and \(q_t\) captures relative signal-quality changes.
\end{enumerate}

\subsection{Why Fixed-LR Sweeps Remain Strong Baselines}

If the goal is only to obtain the best performance in a single fixed setting, sweeping a fixed LR can be effective. A stable fixed LR can avoid high-LR late-stage degradation. However, a fixed-LR sweep does not explain:
\begin{itemize}[leftmargin=2em]
    \item why the same LR can help early training but hurt late training;
    \item whether the selected LR transfers after changing \(N,K\), reward/verifier, model scale, or data distribution;
    \item whether the effective policy-gradient signal changes structurally during training.
\end{itemize}

Thus, the goal of LR scheduling should not be stated as ``fully replacing fixed-LR sweeps.'' A more defensible goal is:
\[
\boxed{
\text{reduce stage-wise LR mismatch, seed sensitivity, and cross-setting tuning cost.}
}
\]

\section{Limitations of Earlier Routes: Batch Split, Momentum, and MARS}

\subsection{Failure of Batch-Split \(r_t\)}

The earlier estimator splits the batch into two halves:
\[
\hat g_A=g_t+\epsilon_A,\qquad
\hat g_B=g_t+\epsilon_B.
\]
It estimates the signal fraction as
\[
\hat r_t^{split}
=
\frac{\hat g_A^\top \hat g_B}
{\frac{1}{2}(\|\hat g_A\|^2+\|\hat g_B\|^2)}.
\]
The numerator expands as
\[
\hat g_A^\top \hat g_B
=
\|g_t\|^2
+
g_t^\top \epsilon_A
+
g_t^\top \epsilon_B
+
\epsilon_A^\top \epsilon_B.
\]
In late-stage RL, the true signal \(\|g_t\|^2\) can become small, so the estimator is dominated by the noise term \(\epsilon_A^\top\epsilon_B\).

Empirically, we observed:
\[
r_t^{split}\approx 0.002\sim 0.004,\qquad
\Pr(g_A^\top g_B>0)\approx 54\%.
\]
This indicates that split-batch alignment is only slightly better than a coin flip and is unsuitable as a step-level LR controller.

\subsection{Why Momentum Should Not Be the Main \(r_t\) Estimator}

Adam momentum
\[
m_t=\beta_1 m_{t-1}+(1-\beta_1)\hat g_t
\]
is useful as an optimizer-side update smoother. However, using \(\|m_t\|^2\) as an estimate of the current \(\|g_t\|^2\) requires a much stronger quasi-static assumption.

Expanding the momentum gives
\[
m_t=(1-\beta_1)\sum_{j=0}^\infty \beta_1^j \hat g_{t-j}.
\]
It tracks the historical average gradient
\[
\bar g_t=(1-\beta_1)\sum_{j=0}^\infty \beta_1^j g_{t-j},
\]
not the current \(g_t\). In RL, the policy distribution, reward structure, and advantage informativeness all change during training. In signal-decay regimes,
\[
\|\bar g_t\|^2>\|g_t\|^2,
\]
so the momentum norm can overestimate current signal power.

Therefore:
\[
\boxed{
\text{momentum is useful as an optimizer or diagnostic, but not as the main } r_t \text{ estimator.}
}
\]

\subsection{Positioning MARS}

MARS is an optimizer-side variance-reduction method that combines scaled stochastic recursive momentum with optimizers such as AdamW, Lion, or Shampoo. It answers:
\[
\boxed{
\text{given a stochastic gradient, how can we obtain a less noisy update direction?}
}
\]
Our current question is:
\[
\boxed{
\text{given the current RL signal quality, how should we modulate the base LR?}
}
\]
These are related but not the same. MARS can serve as an optimizer baseline, but it cannot replace RL noise-source decomposition or LR scheduling.


\section{Bridging the Gap: From Gradient Noise to \(A^2\), \(Q\), and \(H=A^2Q\)}

A key conceptual point is that the abstract gradient-noise term
\[
N_t=\mathbb{E}\|\hat g_t-g_t\|^2
\]
is not directly equal to \(A^2\), \(Q\), or \(A^2Q\). These quantities appear only after expanding the response-level policy-gradient contribution whose finite-sample average forms \(\hat g_t\). The purpose of this section is to make this bridge explicit.

\subsection{Gradient Noise as Finite-Sample Error of Policy-Gradient Contributions}

At the abstract level, the stochastic gradient can be written as
\[
\hat g_t=g_t+\epsilon_t,
\]
where \(g_t=\nabla J(\theta_t)\) is the true policy gradient and \(\epsilon_t\) is the finite-sample estimation error. The noise power in the learning-rate derivation is
\[
N_t
=
\mathbb{E}\|\hat g_t-g_t\|^2
=
\mathbb{E}\|\epsilon_t\|^2.
\]

In policy-gradient training, the batch gradient is an average of response-level contributions. For prompt \(i\) and response \(k\), define
\[
z_{i,k}
=
A_{i,k}\nabla_\theta \log \pi_\theta(y_{i,k}\mid x_i).
\]
Then the batch gradient can be written as
\[
\hat g_t
=
\frac{1}{N}
\sum_{i=1}^{N}
\frac{1}{K}
\sum_{k=1}^{K}
z_{i,k}.
\]
The true policy gradient is
\[
g_t=\mathbb{E}[z].
\]
Therefore, the gradient noise is the finite-sample error of estimating \(\mathbb{E}[z]\) with a finite batch:
\[
\epsilon_t=\hat g_t-g_t.
\]

For a random contribution \(z\), its variance satisfies
\[
\operatorname{Var}(z)
=
\mathbb{E}\|z\|^2
-
\|\mathbb{E}[z]\|^2
=
\mathbb{E}\|z\|^2
-
\|g_t\|^2.
\]
For a batch mean with effective sample size \(n_{\mathrm{eff},t}\), the gradient-noise power roughly scales as
\[
N_t
\approx
\frac{1}{n_{\mathrm{eff},t}}
\left(
\mathbb{E}\|z\|^2-\|g_t\|^2
\right).
\]
When the true signal \(\|g_t\|^2\) is small relative to the second moment, a useful approximation is
\[
N_t
\approx
\frac{1}{n_{\mathrm{eff},t}}
\mathbb{E}\|z\|^2.
\]
Thus, to understand gradient noise, it is natural to inspect the second moment of the response-level policy-gradient contribution.

\subsection{Expanding \(z\) Introduces the Advantage-Energy Term \(A^2\)}

Because
\[
z=A\nabla_\theta\log\pi_\theta(y\mid x),
\]
we have
\[
\|z\|^2
=
A^2
\left\|
\nabla_\theta\log\pi_\theta(y\mid x)
\right\|^2.
\]
Therefore,
\[
\mathbb{E}\|z\|^2
=
\mathbb{E}
\left[
A^2
\left\|
\nabla_\theta\log\pi_\theta(y\mid x)
\right\|^2
\right].
\]

This is the first bridge from abstract gradient noise to RL-specific quantities. The factor \(A^2\) measures the strength of the advantage signal.

In GRPO, if all responses under a prompt receive the same reward, the group-normalized advantages are approximately zero:
\[
A_{i,k}\approx 0.
\]
Such prompt groups provide little or no relative preference signal for the actor update. Therefore, \(A^2\) captures both signal availability and signal concentration:
\begin{enumerate}[leftmargin=2em]
    \item If many responses have \(A^2\approx 0\), the effective policy-gradient signal is sparse.
    \item If only a few responses have very large \(A^2\), the gradient energy becomes heavy-tailed and concentrated.
\end{enumerate}

\subsection{Approximating the Score Norm with a Logit-Space Proxy \(Q\)}

The true parameter-space score norm
\[
\left\|
\nabla_\theta\log\pi_\theta(y\mid x)
\right\|^2
\]
is expensive to compute per response for large language models. We therefore use a logit-space proxy.

For a token with logits \(\ell\), probability vector
\[
p=\operatorname{softmax}(\ell),
\]
and sampled token \(a\), the logit-space score is
\[
\nabla_\ell \log p(a)=e_a-p.
\]
Its squared norm is
\[
q_\tau
=
\|e_a-p\|^2
=
1-2p(a)+\sum_v p(v)^2.
\]
For a full response, define
\[
Q_{i,k}
=
\sum_{\tau=1}^{T_{i,k}}
q_{i,k,\tau}.
\]

Thus, \(Q_{i,k}\) is a score-function / response-length energy proxy. It does not measure whether a response is good or bad; it measures how large the score-side gradient energy is for that response.

Since the true parameter-space score norm is not directly available, \(Q_{i,k}\) is used as a cheap proxy for its relative variation.

\subsection{The Combined Proxy \(H=A^2Q\)}

Combining the advantage-energy factor and the score-energy proxy gives
\[
\|z_{i,k}\|^2
\approx
A_{i,k}^2Q_{i,k}.
\]
We define
\[
\boxed{
H_{i,k}=A_{i,k}^2Q_{i,k}
}
\]
and interpret it as a response-level policy-gradient energy proxy.

Therefore,
\[
\boxed{
A^2 \text{ measures advantage-signal strength,}
}
\]
\[
\boxed{
Q \text{ measures score/length energy,}
}
\]
and
\[
\boxed{
H=A^2Q \text{ approximates response-level policy-gradient energy.}
}
\]

The full bridge can be summarized as
\[
\text{gradient noise}
\rightarrow
\text{finite-sample variance of } z
\rightarrow
\mathbb{E}\|z\|^2
\rightarrow
A^2\|\nabla\log\pi\|^2
\rightarrow
A^2Q.
\]

Thus, \(A^2\), \(Q\), and \(H\) are not themselves the abstract gradient-noise term \(N_t\). They are observable factors that control the second moment of policy-gradient contributions, and therefore help explain finite-sample gradient noise.

\subsection{Connecting \(H\) to Prompt-Level and Rollout-Level Noise}

Once each response has a scalar policy-gradient energy proxy \(H_{i,k}\), we can analyze where its variability comes from. For each prompt, define
\[
\mu_i^H
=
\frac{1}{K}
\sum_{k=1}^{K}
H_{i,k}.
\]
The across-prompt component is
\[
V_{\mathrm{between}}^H
=
\operatorname{Var}_i(\mu_i^H).
\]
This measures how different prompts vary in their average policy-gradient energy proxy.

The within-prompt rollout component is
\[
V_{\mathrm{within,raw}}^H
=
\mathbb{E}_i
\left[
\operatorname{Var}_k(H_{i,k})
\right].
\]
Since the prompt-level estimator averages over \(K\) responses, the batch-level within contribution should be compared after the \(1/K\) correction:
\[
V_{\mathrm{within,corr}}^H
=
\frac{V_{\mathrm{within,raw}}^H}{K}.
\]
The corrected between fraction is
\[
f_{\mathrm{between}}^H
=
\frac{
V_{\mathrm{between}}^H
}{
V_{\mathrm{between}}^H+V_{\mathrm{within,raw}}^H/K+\epsilon
}.
\]

This decomposition does not claim to exactly equal the full parameter-space gradient covariance. Instead, it is a scalar policy-gradient energy proxy decomposition. It is useful for diagnosing whether the effective gradient energy is dominated by:
\begin{enumerate}[leftmargin=2em]
    \item prompt-to-prompt heterogeneity;
    \item within-prompt rollout variability;
    \item advantage-signal sparsity or heavy tails;
    \item score/length amplification.
\end{enumerate}

In later empirical sections, we analyze \(A^2\), \(Q\), and \(H=A^2Q\) separately to determine which factor drives the observed policy-gradient energy concentration.


\section{Noise Decomposition in RL: Two Levels}

\subsection{Level 1: Sampling Axis}

In GRPO, a prompt \(x_i\) has \(K\) responses:
\[
y_{i,1},\ldots,y_{i,K}.
\]
The response-level policy-gradient contribution is
\[
z_{i,k}=A_{i,k}\nabla_\theta \log \pi_\theta(y_{i,k}|x_i).
\]
The prompt-level contribution is
\[
U_i=\frac{1}{K}\sum_{k=1}^K z_{i,k}.
\]
The batch gradient is
\[
\hat g=\frac{1}{N}\sum_{i=1}^N U_i.
\]
Define
\[
\mu(x)=\E[U|x].
\]
By the law of total variance,
\[
\Var(U)
=
\underbrace{\Var_x(\mu(x))}_{\text{between-prompt noise}}
+
\underbrace{\E_x[\Var(U|x)]}_{\text{within-prompt rollout noise}}.
\]
Under an approximate independence assumption across responses,
\[
N_t
\approx
\frac{1}{N}\sigma^2_{\mathrm{prompt},t}
+
\frac{1}{NK}\sigma^2_{\mathrm{rollout},t}.
\]

Thus the two top-level sampling noise sources are:
\[
\boxed{\text{between-prompt sampling noise}}
\]
and
\[
\boxed{\text{within-prompt rollout sampling noise}}.
\]

\subsection{Level 2: Factor Axis}

The true response-level gradient energy is
\[
\|z_{i,k}\|^2
=
A_{i,k}^2
\left\|
\nabla_\theta\log\pi_\theta(y_{i,k}|x_i)
\right\|^2.
\]
The parameter-space score norm is expensive, so we use a logit-space score/length proxy \(Q_{i,k}\).

For a token \(\tau\), let the probability vector be \(p\) and the sampled token be \(a\). The logit-space score is
\[
\nabla_\ell \log p(a)=e_a-p,
\]
with squared norm
\[
q_\tau
=
\|e_a-p\|^2
=
1-2p(a)+\sum_v p(v)^2.
\]
The response-level score energy is
\[
Q_{i,k}=\sum_{\tau=1}^{T_{i,k}} q_{i,k,\tau}.
\]
We define
\[
\boxed{
H_{i,k}=A_{i,k}^2Q_{i,k}.
}
\]
Here:
\begin{itemize}[leftmargin=2em]
    \item \(A_{i,k}\): GRPO advantage;
    \item \(A_{i,k}^2\): advantage energy, or reward/advantage signal strength;
    \item \(Q_{i,k}\): score-function / response-length energy proxy;
    \item \(H_{i,k}=A_{i,k}^2Q_{i,k}\): policy-gradient energy proxy.
\end{itemize}

The second-level question is:
\[
\boxed{
\text{is the noise driven by advantage signal, score/length energy, or their interaction?}
}
\]

\section{Stage-2 Setup and Overall Performance}

\subsection{Experimental Setup}

The total number of responses per batch is fixed:
\[
N\times K=1024.
\]
We compare three configurations:

\begin{table}[htbp]
\centering
\small
\begin{tabular}{@{}lccc>{\raggedright\arraybackslash}p{0.42\linewidth}@{}}
\toprule
Config & \(N\) & \(K\) & MB & Purpose \\
\midrule
n256k4 & 256 & 4  & 64 & High prompt diversity; low within-prompt averaging.\\
n128k8 & 128 & 8  & 32 & Baseline configuration.\\
n64k16 & 64  & 16 & 16 & Low prompt diversity; high within-prompt averaging.\\
\bottomrule
\end{tabular}
\caption{Stage-2 \(N\times K\) rerun setup. Here MB denotes the mini-batch size.}
\end{table}

Key setup:
\begin{itemize}[leftmargin=2em]
    \item Model: DeepSeek-R1-Distill-Qwen-1.5B;
    \item LR: constant \(3.1\times10^{-6}\), a stable LR regime;
    \item usable runs: 9 runs, 3 seeds per config;
    \item target length: roughly 300 updates;
    \item validation: core avg5 = mean(AIME, AIME2025, GPQA, MINERVA, OLYMPIAD\_BENCH) acc/mean@16.
\end{itemize}

\subsection{Validation Performance}

\begin{table}[h]
\centering
\begin{tabular}{lrrrrr}
\toprule
Config & Seeds & Mean Best Avg5 & Mean Final Avg5 & Mean Drop & Std Final\\
\midrule
n256k4 & 3 & 0.3290 & 0.3282 & -0.0008 & 0.0021\\
n128k8 & 3 & \textbf{0.3341} & \textbf{0.3315} & -0.0026 & \textbf{0.0016}\\
n64k16 & 3 & 0.3312 & 0.3279 & -0.0033 & 0.0091\\
\bottomrule
\end{tabular}
\caption{Validation performance across \(N,K\) configurations.}
\end{table}

The mean final differences across the three configurations are small, around \(0.003\). The n64k16 configuration has the largest seed-level standard deviation. Thus, one should not claim a statistically strong winner. Numerically, n128k8 is slightly better and appears to be a reasonable tradeoff under the current budget.

\section{Sampling Axis: Between/Within Prompt Decomposition}

\subsection{Corrected Decomposition for \(H=A^2Q\)}

For any \(X\in\{A^2,Q,H\}\), define the prompt-level mean:
\[
\mu_i^X=\frac{1}{K}\sum_{k=1}^K X_{i,k}.
\]
The between component is
\[
V_{\mathrm{between}}^X=\Var_i(\mu_i^X).
\]
The raw within component is
\[
V_{\mathrm{within,raw}}^X=\E_i[\Var_k(X_{i,k})].
\]
For batch-level comparison, the within term should be corrected as
\[
V_{\mathrm{within,corr}}^X=\frac{V_{\mathrm{within,raw}}^X}{K}.
\]
Thus,
\[
f_{\mathrm{between}}^X
=
\frac{V_{\mathrm{between}}^X}
{V_{\mathrm{between}}^X+V_{\mathrm{within,raw}}^X/K+\epsilon}.
\]

For \(H=A^2Q\), the empirical results are:

\begin{table}[h]
\centering
\begin{tabular}{lrrrrr}
\toprule
Config & \(K\) & Early & Mid & Late & Overall\\
\midrule
n256k4 & 4 & 0.642 & 0.673 & 0.681 & \textbf{0.664}\\
n128k8 & 8 & 0.527 & 0.556 & 0.572 & 0.550\\
n64k16 & 16 & 0.487 & 0.496 & 0.514 & \textbf{0.499}\\
\bottomrule
\end{tabular}
\caption{Batch-corrected between-prompt fraction for \(H=A^2Q\).}
\end{table}

The key conclusion is:
\[
\boxed{
K=4 \text{ is prompt-noise dominated, while } K=8/16 \text{ are roughly balanced.}
}
\]
Therefore, the issue is not that within-prompt rollout noise always dominates. Instead, \(N/K\) controls the tradeoff between prompt diversity and within-prompt comparison quality.

\subsection{Between Fraction by Factor}

\begin{table}[h]
\centering
\begin{tabular}{lrrrl}
\toprule
Factor & n256k4 & n128k8 & n64k16 & Interpretation\\
\midrule
\(A^2\) & 0.677 & 0.528 & 0.430 & most sensitive to \(K\)\\
\(Q\) & 0.914 & 0.947 & 0.968 & mostly a prompt-level property\\
\(H=A^2Q\) & 0.664 & 0.550 & 0.499 & combined result\\
\bottomrule
\end{tabular}
\caption{Between-prompt fractions for \(A^2,Q,H\).}
\end{table}

This suggests:
\[
\boxed{
Q \text{ is strongly prompt-structured, but its independent variation is not the main driver of } H.
}
\]

\section{Factor Axis: Independent Analysis of \(A^2\), \(Q\), and \(H\)}

\subsection{CV\(^2\): Relative Variation}

\begin{table}[h]
\centering
\begin{tabular}{lrrr}
\toprule
Metric & n256k4 & n128k8 & n64k16\\
\midrule
\(\CV^2(A^2)\) & 3.17 & 4.59 & \textbf{6.84}\\
\(\CV^2(Q)\) & 0.34 & 0.36 & 0.36\\
\(\CV^2(H)\) & 4.72 & 7.08 & \textbf{10.76}\\
\bottomrule
\end{tabular}
\caption{Relative variation of \(A^2,Q,H\).}
\end{table}

The relative variation of \(A^2\) is much larger than that of \(Q\):
\[
\CV^2(A^2)\gg \CV^2(Q).
\]
Thus, differences in response-level gradient energy are mainly driven by the advantage side rather than the score/length side.

\subsection{Counterfactual Attribution}

We decompose the variance of
\[
H=A^2Q
\]
into three diagnostic components:
\begin{itemize}[leftmargin=2em]
    \item \(A^2\)-only: fix \(Q\) to its mean and allow \(A^2\) to vary;
    \item \(Q\)-only: fix \(A^2\) to its mean and allow \(Q\) to vary;
    \item interaction: the remaining multiplicative coupling.
\end{itemize}

\begin{table}[h]
\centering
\begin{tabular}{lrrr}
\toprule
Attribution & n256k4 & n128k8 & n64k16\\
\midrule
\(A^2\)-only frac & 0.531 & 0.499 & 0.488\\
\(Q\)-only frac & 0.062 & 0.041 & 0.028\\
Interaction frac & 0.407 & 0.460 & 0.484\\
\bottomrule
\end{tabular}
\caption{Counterfactual attribution of \(H=A^2Q\) variance.}
\end{table}

The conclusion is:
\[
\boxed{
\text{about half of }H\text{ variance comes from } A^2,\quad
Q\text{-only explains only }3\%\sim6\%,\quad
\text{and interaction explains }40\%\sim48\%.
}
\]
Therefore, \(Q\) mainly acts as an amplifier through multiplication with \(A^2\), rather than as an independent primary driver.

\subsection{Training-Time Decay}

\begin{table}[h]
\centering
\begin{tabular}{llrrr}
\toprule
Component & Metric & n256k4 & n128k8 & n64k16\\
\midrule
\(A^2\) & \(\E[A^2]\) decay & -25.7\% & -23.0\% & -17.2\%\\
\(Q\) & score energy decay & -21.3\% & -27.0\% & -13.5\%\\
\(H\) & \(\E[H]\) decay & -49.8\% & -48.7\% & -35.0\%\\
\bottomrule
\end{tabular}
\caption{Early-to-late decay of \(A^2,Q,H\).}
\end{table}

There is no simple pattern of ``score energy stays high while signal weakens.'' Instead, both \(A^2\) and \(Q\) decay, and their product \(H\) decays even more. However, the distribution of \(A^2\) becomes more heavy-tailed, causing the remaining effective signal to become increasingly concentrated.

\section{Reward Informativeness: Signal Sparsity}

\begin{table}[h]
\centering
\begin{tabular}{lrrr}
\toprule
Metric & n256k4 & n128k8 & n64k16\\
\midrule
frac\_informative & 0.401 & 0.536 & \textbf{0.648}\\
frac\_all\_wrong & 0.349 & 0.279 & 0.229\\
frac\_all\_correct & 0.249 & 0.186 & 0.123\\
reward\_n\_unique/mean & 1.50 & 1.86 & 2.29\\
reward\_std/median & 0.054 & 0.533 & \textbf{0.597}\\
\bottomrule
\end{tabular}
\caption{Reward informativeness statistics.}
\end{table}

The key finding is:
\[
\boxed{
K \uparrow \Rightarrow \text{more prompt groups have mixed rewards and richer effective advantage signal.}
}
\]
However, higher reward informativeness does not necessarily imply better validation performance, because larger \(K\) reduces the number of prompts \(N\) and increases the heavy-tailed concentration of \(A^2/H\).

\subsection{Training-Time Degradation}

The fraction of informative prompt groups decreases during training in all configurations:
\[
\text{n256k4: } -8.5\%,\qquad
\text{n128k8: } -9.7\%,\qquad
\text{n64k16: } -4.7\%.
\]
This suggests that as the model improves, more prompt groups become all-correct or all-wrong, reducing effective group-level relative signal.

\section{Gradient-Energy Concentration: Effective Batch Size Is Much Smaller Than Nominal Batch Size}

\subsection{Response-Level Concentration}

The top-10\% responses contribute most of \(|H|\):

\begin{table}[h]
\centering
\begin{tabular}{lrrrrrr}
\toprule
Config & step 10 & 50 & 100 & 150 & 200 & 250\\
\midrule
n256k4 & 0.485 & 0.637 & 0.659 & 0.617 & 0.664 & 0.658\\
n128k8 & 0.583 & 0.695 & 0.654 & 0.680 & 0.715 & 0.699\\
n64k16 & \textbf{0.640} & \textbf{0.741} & \textbf{0.771} & \textbf{0.724} & \textbf{0.724} & \textbf{0.714}\\
\bottomrule
\end{tabular}
\caption{Top-10\% response share of total \(|H|\).}
\end{table}

For \(K=16\), only 10\% of responses contribute around \(71\%-77\%\) of \(|H|\).

\subsection{Prompt-Level Concentration}

At step 150:

\begin{table}[h]
\centering
\begin{tabular}{lrrr}
\toprule
Config & prompts for 50\% within-H-var & prompts for 80\% & informative/total\\
\midrule
n256k4 & 21/256 (8.2\%) & 48/256 (18.8\%) & 115/256 (44.9\%)\\
n128k8 & 9/128 (7.0\%) & 21/128 (16.4\%) & 66/128 (51.6\%)\\
n64k16 & 5/64 (7.8\%) & 13/64 (20.3\%) & 40/64 (62.5\%)\\
\bottomrule
\end{tabular}
\caption{Prompt concentration of within-H variance at step 150.}
\end{table}

Only around \(7\%-8\%\) of prompts contribute 50\% of within-H variance. This implies:
\[
\boxed{
\text{the nominal batch size is large, but the effective gradient-contributing sample size is much smaller.}
}
\]

\section{Summary of the Identified Problem}

The problem we have identified has three layers.

\subsection{Layer 1: Sampling Tradeoff}

Under a fixed rollout budget:
\[
N\times K=1024.
\]
Increasing \(K\) has two opposite effects:
\begin{itemize}[leftmargin=2em]
    \item it improves within-prompt comparison quality and increases frac\_informative;
    \item it reduces prompt diversity and increases \(A^2/H\) heavy-tail concentration risk.
\end{itemize}
Therefore, larger \(K\) is not always better. Under the current stable-LR setting, \(N=128,K=8\) appears to be a reasonable empirical tradeoff.

\subsection{Layer 2: Advantage Signal Is the Main Driver}

The relative variation of \(A^2\) is much larger than that of \(Q\), and \(A^2\)-only explains about half of the variance of \(H\). The main issue is therefore on the reward/advantage side:
\[
\boxed{
\text{effective advantage signal is sparse and heavy-tailed.}
}
\]

\subsection{Layer 3: Gradient-Energy Concentration}

A small number of prompts/responses dominate most of \(H\), causing the effective sample size to be much smaller than the nominal batch size. A fixed LR can become overly sensitive to such heavy-tailed batches.

\section{Sampling vs LR: How Should We Position the Intervention?}

\subsection{Sampling Is the Root-Cause Intervention}

The experimental evidence points strongly toward sampling structure:
\[
\boxed{
\text{sampling changes the effective signal and noise themselves.}
}
\]
Potential sampling-side methods include:
\begin{itemize}[leftmargin=2em]
    \item dynamic \(K\);
    \item prompt difficulty balancing;
    \item historical success-rate sampling;
    \item informative prompt sampling;
    \item heavy-tail prompt downweighting;
    \item response-level clipping or rejection.
\end{itemize}

\subsection{LR Is a Risk-Control Intervention}

LR does not change the signal quality of the current batch. It only reduces the step size when the current batch is unreliable:
\[
\boxed{
\text{LR is a low-intrusion risk controller, not a root-cause solution.}
}
\]
Therefore, an LR method is valuable only if it can reduce high-LR degradation, peak-to-final drop, or seed variance without changing the rollout pipeline.

\subsection{Recommended Experiment: Orthogonal Comparison}

To decide whether sampling or LR is the stronger intervention, we should run:

\[
\begin{array}{c|c}
\text{Sampling} & \text{LR}\\
\hline
\text{baseline sampling} & \text{fixed LR}\\
\text{improved sampling} & \text{fixed LR}\\
\text{baseline sampling} & \text{adaptive LR}\\
\text{improved sampling} & \text{adaptive LR}
\end{array}
\]

This answers:
\begin{enumerate}[leftmargin=2em]
    \item whether sampling alone is sufficient;
    \item whether LR has independent benefit;
    \item whether sampling and LR are complementary;
    \item whether the main algorithmic contribution should focus on sampling or LR.
\end{enumerate}

\section{A New LR Direction: Signal-Quality-Aware LR Scaling}

\subsection{Basic Form}

We no longer attempt to estimate an exact \(r_t\). Instead, we use:
\[
\boxed{
\alpha_t=\alpha_{\mathrm{base}} q_t.
}
\]
Here \(\alpha_{\mathrm{base}}\) absorbs the average \(1/L_t\) scale, while \(q_t\) is a signal-quality gate.

\subsection{Signal Availability Gate}

Define signal availability:
\[
I_t\in\left\{
\mathrm{reward\_std/median},\ 
\mathrm{frac\_informative},\ 
\E[A^2]
\right\}.
\]
Then:
\[
q_t^{info}
=
\clip\left(
\frac{\mathrm{EMA}(I_t)}
{I_{\mathrm{ref}}},
q_{\min}^{info},
1
\right).
\]

\subsection{Concentration Penalty}

Define concentration risk:
\[
C_t\in\left\{
\CV^2(A^2),\ 
\CV^2(H),\ 
\text{top-10\% }H\text{ share}
\right\}.
\]
Then:
\[
q_t^{conc}
=
\clip\left[
\left(
\frac{C_{\mathrm{ref}}}{\mathrm{EMA}(C_t)+\epsilon}
\right)^\gamma,
q_{\min}^{conc},
1
\right].
\]

The combined LR is:
\[
\boxed{
\alpha_t
=
\alpha_{\mathrm{base}}\,
q_t^{info}\,
q_t^{conc}.
}
\]

Intuition:
\begin{itemize}[leftmargin=2em]
    \item reduce LR when effective signal is insufficient;
    \item reduce LR when gradient energy is overly concentrated.
\end{itemize}

\section{Minimal Algorithmic Experiment Plan}

\subsection{Stage A: Logging-Only}

On existing fixed-LR runs, compute offline:
\[
q_t^{info},\qquad q_t^{conc},\qquad q_t=q_t^{info}q_t^{conc}.
\]
Check whether \(q_t\) decreases before validation drop, entropy collapse, or ratio-tail expansion.

\subsection{Stage B: Info-Only LR Gate}

Use an aggressive base LR:
\[
\alpha_{\mathrm{base}}=10^{-5},
\]
and set
\[
\alpha_t=\alpha_{\mathrm{base}}q_t^{info}.
\]
Test whether it preserves early fast learning while reducing late-stage regression.

\subsection{Stage C: Info + Concentration LR Gate}

Use
\[
\alpha_t=\alpha_{\mathrm{base}}q_t^{info}q_t^{conc}.
\]
Compare peak-to-final drop, seed variance, KL/ratio tail, and final validation.

\subsection{Stage D: Comparison Against Fixed-LR Sweep}

Compare:
\begin{enumerate}[leftmargin=2em]
    \item stable fixed LR;
    \item aggressive fixed LR;
    \item oracle fixed LR sweep;
    \item aggressive LR + info gate;
    \item aggressive LR + info + concentration gate.
\end{enumerate}

Evaluation should include:
\[
\text{peak-to-final drop},\quad
\text{seed variance},\quad
\text{late-stage degradation},\quad
\text{cross-setting transfer}.
\]

\section{Related Work Positioning}

\subsection{Gradient Noise Scale}

Gradient-noise-scale and critical-batch-size work connects batch size, noise scale, and training efficiency. However, it mostly concerns general stochastic optimization and does not explicitly model GRPO prompt/rollout/advantage structure.

\subsection{Policy-Gradient SNR and Adaptive Step Size}

Prior work on policy-gradient SNR and adaptive step size shows that policy-gradient step size should depend on performance bounds and gradient uncertainty. This supports our \(L\)-smoothness and signal-fraction motivation.

\subsection{TRPO/PPO}

TRPO and PPO control policy-space update size using KL trust regions or clipping. Our work is different: we modulate the optimizer LR using RL-specific signal-quality diagnostics.

\subsection{MARS and Optimizer-Side Variance Reduction}

MARS is an optimizer-side variance-reduction method. It can serve as an AdamW baseline, but it is not an LR scheduler and cannot replace RL noise-source decomposition.

\subsection{GRPO/DAPO and Dynamic Sampling}

GRPO constructs advantages using group-relative rewards. DAPO and related LLM RL recipes emphasize dynamic sampling, clipping strategies, entropy, and diversity. This is consistent with our finding that sampling structure and advantage informativeness are central to stable LLM RL.

\section{Conclusion}

We have reframed the problem from:
\[
\text{``How do we estimate a generic } r_t \text{?''}
\]

to the following question:

\begin{center}
\fbox{%
\begin{minipage}{0.88\linewidth}
\centering
How do we detect sparsity and concentration in the effective GRPO
policy-gradient signal, and decide whether to intervene through
sampling or learning-rate control?
\end{minipage}
}
\end{center}

The main empirical findings are:
\begin{enumerate}[leftmargin=2em]
    \item batch-split \(r_t\) is too noisy to serve as the main LR controller;
    \item \(N/K\) controls the tradeoff between prompt diversity and within-prompt comparison quality;
    \item the relative variation of \(A^2\) is much larger than that of \(Q\), so the advantage side is the main driver;
    \item \(H=A^2Q\) is highly concentrated on a small subset of responses and prompts;
    \item effective batch size is much smaller than nominal batch size;
    \item sampling-side correction is the root-cause intervention, while LR-side correction is a low-intrusion risk-control mechanism.
\end{enumerate}

Thus, the earlier PG-Moment/\(\lambda_M\) calibration route should not be the main algorithmic direction. The next step should be an orthogonal comparison between:
\[
\boxed{
\text{sampling improvement}
}
\quad \text{and}\quad
\boxed{
\text{signal-quality-aware LR scaling}.
}
\]

If we continue with LR, the recommended first method is:
\[
\boxed{
\alpha_t
=
\alpha_{\mathrm{base}}\,
q_t^{info}\,
q_t^{conc},
}
\]
where \(q_t^{info}\) responds to effective advantage-signal availability and \(q_t^{conc}\) penalizes gradient-energy concentration risk.

\newpage
\begin{thebibliography}{9}

\bibitem{mccandlish2018}
S. McCandlish et al.
\newblock An Empirical Model of Large-Batch Training.
\newblock arXiv:1812.06162, 2018.

\bibitem{trpo}
J. Schulman et al.
\newblock Trust Region Policy Optimization.
\newblock arXiv:1502.05477, 2015.

\bibitem{ppo}
J. Schulman et al.
\newblock Proximal Policy Optimization Algorithms.
\newblock arXiv:1707.06347, 2017.

\bibitem{svrpg}
C. Papini et al.
\newblock Stochastic Variance-Reduced Policy Gradient.
\newblock arXiv:1806.05618, 2018.

\bibitem{deepseekmath}
Z. Shao et al.
\newblock DeepSeekMath: Pushing the Limits of Mathematical Reasoning in Open Language Models.
\newblock arXiv:2402.03300, 2024.

\bibitem{mars}
H. Yuan et al.
\newblock MARS: Unleashing the Power of Variance Reduction for Training Large Models.
\newblock arXiv:2411.10438, 2024.

\bibitem{dapo}
Q. Yu et al.
\newblock DAPO: An Open-Source LLM Reinforcement Learning System at Scale.
\newblock arXiv:2503.14476, 2025.

\end{thebibliography}

\end{document}
